\chapter{Collections}

wick has two container types, the list and the map, and one idiom --- parallel
lists --- that stands in for the structs it does not yet have. This chapter
covers all three, along with the iteration patterns that go with 0-based
indexing and the deliberate choice to make every fallible lookup an optional.

\section{Lists}

A list is a typed, ordered, 0-based sequence. Its element type is one of the
three scalar types; a list literal infers its element type from its contents,
and an empty literal must be annotated.

\begin{lstlisting}[language=wick]
let xs = [1, 2, 3]            // list<num>, element type inferred
let names = ["a", "b"]        // list<str>
let flags: list<bool> = []    // empty literal REQUIRES the annotation
\end{lstlisting}

\noindent
The annotation on the empty list is not optional politeness; there is nothing
in `[]` to infer from, and the compiler says so:
\texttt{empty [] needs a type: let xs: list<num> = []}. A literal with mixed
element types does not compile either --- lists are homogeneous ---
and reports \texttt{list elements must share one type}, so `[1, "a"]` is
rejected.

\paragraph{Operating on lists.} The core operations are built-ins, dispatched
by the static element type:

\begin{lstlisting}[language=wick]
push(xs, 4)                  // append; type-checked against the element type
let last = pop(xs) ?? 0      // remove and return the last; T? (nil if empty)
xs[0] = 10                   // index assignment
let head = xs[0]             // index read
let n = len(xs)              // number of elements
\end{lstlisting}

\noindent
`push` appends a value that must match the element type --- `push(flags, 3)`
does not compile. `pop` removes and returns the last element, but a list may
be empty, so its result is an optional `T?`, `nil` when there is nothing to
pop; you handle it like any optional, usually with `?? default`. `len` returns
the element count as a `num`.

\paragraph{Indexing is checked at run time.} Reading or writing an index
outside \texttt{[0, len)} is a \emph{runtime} error, not undefined behaviour
and not a silent `nil`. It stops the frame and shows on the error screen with
the line number and the message \texttt{list index N out of range (len M)}.
The static type system cannot know the value of an index in general, so this
one check remains dynamic --- but it is a check, with a precise message, not a
memory-safety hole.

\paragraph{Element types.} List elements are `num`, `bool`, or `str`. There
are no lists of lists in v0.1: `list<list<num>>` reports
\texttt{container elements must be num, bool, or str}. The workaround, when
you genuinely need a grid or a jagged structure, is index arithmetic on a flat
list, shown at the end of this chapter.

\section{Maps}

A map is a string-keyed, value-typed collection. Keys are always `str`;
values are one of the three scalar types.

\begin{lstlisting}[language=wick]
let scores = ["ana": 3, "bo": 5]   // map<num>
scores["cid"] = 1                  // insert or update
let s = scores["dee"] ?? 0         // get is ALWAYS T? -- missing key is nil
let n = len(scores)                // number of entries
\end{lstlisting}

\noindent
A map literal pairs string-literal keys with values; the keys in a literal
must be string literals, not computed expressions
(\texttt{map keys must be str literals}). To insert under a computed key, use
index assignment: `m[k] = v`. The values in a literal must share one type,
just as list elements do.

\paragraph{Lookup is always optional.} The single most important fact about
maps is that `m[key]` has type `T?`, always, even for a key you just inserted.
The type system cannot prove a key is present, so it treats every lookup as
possibly missing and hands you a `T?`. This is deliberate: the missing-key
case --- the one everybody forgets --- is forced into view at the call site,
where `??` handles it in three characters. There is no ``it was there a moment
ago'' surprise, because the type never let you assume it.

\begin{lstlisting}[language=wick]
scores["ana"] = 10
let a = scores["ana"] ?? 0    // still T?; the ?? is required, not optional
\end{lstlisting}

\section{Records via parallel lists}

wick v0.1 has no structs. When you would reach for a record --- a lamp with a
position and a lit flag, an enemy with health and a state --- you use
\emph{parallel lists}: several lists of scalars, kept the same length, where a
single index names one logical object across all of them.

\begin{lstlisting}[language=wick]
let lamp_x: list<num> = []
let lamp_z: list<num> = []
let lamp_lit: list<bool> = []
// index i is one lamp: its x is lamp_x[i], its z is lamp_z[i],
// its lit flag is lamp_lit[i]
push(lamp_x, 3.4)  push(lamp_z, 3.4)  push(lamp_lit, true)
push(lamp_x, -3.4) push(lamp_z, 3.4)  push(lamp_lit, true)
\end{lstlisting}

\noindent
To operate on ``a lamp'' you operate on one index across the lists:

\begin{lstlisting}[language=wick]
fn relight(i: num) {
  lamp_lit[i] = true
}

fn draw_lamps() {
  for i in 0..len(lamp_x) {
    if lamp_lit[i] {
      lt.draw(cube, lamp_x[i], 0.62, lamp_z[i], 0, 0, 0,
              0.14, 0.2, 0.14, 1.0, 0.8, 0.3)
    }
  }
}
\end{lstlisting}

\noindent
The honest assessment, which Chapter~8 makes at length: this reads less well
than Lua's list of tables, where a lamp is one value with named fields. It is
the top of wick's roadmap. But it is fully typed --- `lamp_lit` is a
`list<bool>`, so `lamp_lit[i]` is a `bool` the compiler checks --- and for the
scale of a game script it is workable. The one rule you must keep yourself is
that the lists stay the same length; a `push` to one is a `push` to all.

\section{Iteration patterns}

There is one loop form for collections: iterate the indices, and index in.

\begin{lstlisting}[language=wick]
for i in 0..len(xs) {
  use(xs[i])
}
\end{lstlisting}

\noindent
Because the range is end-exclusive and lists are 0-based, `0..len(xs)` visits
exactly the valid indices, with no off-by-one. There is no `for x in list` in
v0.1; the index form is the whole story, and it has the advantage that the
index `i` is in hand when you need it --- for parallel lists you need it always,
and for finding-the-nearest problems you need it to return.

\paragraph{Search that returns an index.} A common shape is ``find the element
matching some criterion, or report that there is none.'' With no multiple
returns and no `for x in list`, you write a loop that carries a best-so-far
index and a sentinel for ``none'':

\begin{lstlisting}[language=wick]
fn nearest_unlit(): num {        // returns a lamp index, or -1
  let bi = 0 - 1
  let bd = 0.0
  for i in 0..len(lamp_x) {
    if not lamp_lit[i] {
      let dx = lamp_x[i] - px
      let dz = lamp_z[i] - pz
      let d = dx * dx + dz * dz
      if bi < 0 or d < bd {
        bi = i
        bd = d
      }
    }
  }
  return bi
}
\end{lstlisting}

\noindent
The caller checks the sentinel before using the result:

\begin{lstlisting}[language=wick]
let i = nearest_unlit()
if i >= 0 {
  // i is a valid index
}
\end{lstlisting}

\noindent
Note that `bi` is written `0 - 1` rather than `-1` in the initialiser; either
is legal --- unary minus exists --- and the game code happens to spell it as a
subtraction. The `-1` sentinel is a plain `num`, not an optional, because it
is cheaper to test `i >= 0` than to thread a `num?` through the caller; this
is the idiomatic stand-in for the multiple-return the language lacks.

\section{Nested data without nested containers}

When you truly need a two-dimensional structure --- a tile grid, say --- and
`list<list<num>>` is unavailable, flatten it: store row-major in a single
list and compute the index.

\begin{lstlisting}[language=wick]
let W = 16
let H = 12
let tiles: list<num> = []
for i in 0..W * H { push(tiles, 0) }

fn get(tx: num, ty: num): num {
  return tiles[ty * W + tx]
}

fn set(tx: num, ty: num, v: num) {
  tiles[ty * W + tx] = v
}
\end{lstlisting}

\noindent
This is the same trick you would use in a systems language over a raw buffer,
and it composes fine with the runtime bounds check: an out-of-range
`ty * W + tx` is caught exactly as any other out-of-range index. It is more
verbose than a nested container would be, and it is the honest v0.1 answer
until nested containers arrive.
